\documentclass[12pt]{article}
% Préambule commun à tous les documents extraits de « La Revue CAPES »
\usepackage[T1]{fontenc}
\usepackage[utf8]{inputenc}
\usepackage{lmodern}
\usepackage{amsmath,amssymb,amsthm,amsfonts,mathrsfs}
\usepackage{setspace}
\usepackage{listings}
\usepackage{pgfplots}
\pgfplotsset{compat=1.18}
\usepackage{longtable}
\usepackage{adjustbox}
\usepackage{lettrine}
\usepackage{tkz-tab}
\usepackage{changepage}
\usepackage{pdflscape}
\usepackage{graphicx}
\usepackage{tikz}
\usepackage{xcolor}
\usepackage[a4paper,top=2cm,bottom=2.5cm,left=2.5cm,right=2cm]{geometry}
\usepackage{fancyhdr}
\usepackage{draftwatermark}
\SetWatermarkText{La RevueCapes}
\SetWatermarkLightness{0.9}
\SetWatermarkScale{2}
\usepackage{hyperref}
\hypersetup{colorlinks=true,linkcolor=blue!50!black,urlcolor=blue!50!black}

\definecolor{revue}{RGB}{20,60,120}

\pagestyle{fancy}
\fancyhf{}
\renewcommand{\headrulewidth}{0.4pt}
\setlength{\headheight}{15pt}

% \entete{Matière}{Titre}{Type de document}
\newcommand{\entete}[3]{%
  \fancyhead[L]{\small\textsc{La Revue CAPES} -- #1}%
  \fancyhead[R]{\small #2}%
  \fancyfoot[C]{\thepage}%
  \thispagestyle{plain}%
  \begin{center}
    {\color{revue}\rule{\textwidth}{1.2pt}}\\[4pt]
    {\small\textsc{Concours directs CAP-CEG / CAPES -- Burkina Faso}}\\[6pt]
    {\LARGE\bfseries\color{revue} #2}\\[6pt]
    {\large\itshape #1 -- #3}\\[2pt]
    {\color{revue}\rule{\textwidth}{1.2pt}}
  \end{center}
  \vspace{0.5cm}%
}

% Encadré pour les remarques ajoutées dans les corrigés
\newcommand{\remarque}[1]{\par\medskip\noindent\fcolorbox{revue}{revue!6}{\parbox{\dimexpr\linewidth-2\fboxsep-2\fboxrule}{\textbf{\color{revue}Remarque.} #1}}\par\medskip}


\begin{document}
\entete{Culture générale}{Corrigé du sujet type de dissertation n°13}{Correction}
\begin{spacing}{1.5}

\textbf{Introduction}

Depuis plusieurs décennies, la présence militaire française au Sahel, notamment à travers l’opération Barkhane, symbolisait l’engagement international dans la lutte contre le terrorisme. Cependant, face à des résultats mitigés et une montée de l’hostilité populaire envers cette présence, plusieurs États sahéliens, dont le Burkina Faso et le Mali, ont choisi de rompre leurs alliances traditionnelles avec la France pour se tourner vers la Russie, notamment à travers le groupe paramilitaire Wagner. Ce repositionnement stratégique soulève des questions fondamentales sur ses implications pour la stabilité régionale et la lutte contre le terrorisme. Quels en sont les impacts sur une région déjà fragilisée par des crises multiples ?

\quad

\textbf{Impacts sur la stabilité régionale}

Le retrait de la présence militaire française et le recours à des forces telles que le groupe Wagner reconfigurent les équilibres géopolitiques au Sahel. Tout d’abord, cette transition a entraîné une détérioration des relations diplomatiques entre les États sahéliens concernés et leurs anciens partenaires occidentaux, notamment la France et l’Union européenne. Cette rupture pourrait limiter l’accès aux financements et aux soutiens techniques nécessaires pour stabiliser la région.

Par ailleurs, l’implication du groupe Wagner, connu pour ses pratiques controversées dans d’autres zones de conflit, comme en République centrafricaine, suscite des craintes quant à la montée des violations des droits humains. Ces abus pourraient exacerber les tensions entre les populations locales et leurs gouvernements, alimentant ainsi un cycle de méfiance et d’instabilité.

Sur le plan régional, le rapprochement avec la Russie pourrait renforcer les divisions entre les États sahéliens, certains privilégiant toujours des partenariats occidentaux. Cela risque d'affaiblir les mécanismes de coopération régionale nécessaires pour répondre aux défis transnationaux, notamment le terrorisme.

\quad

\textbf{Conséquences sur la lutte contre le terrorisme}

En ce qui concerne la lutte contre le terrorisme, les résultats de la coopération avec le groupe Wagner restent incertains. Si ce dernier prétend offrir des solutions militaires rapides et agressives, ses interventions ailleurs montrent une approche principalement axée sur la répression violente. Cette stratégie, sans accompagnement d’initiatives de développement et de réconciliation, risque de radicaliser davantage les populations locales, renforçant ainsi les groupes terroristes.

De plus, l’absence de coordination entre les nouveaux partenaires (Russie et Wagner) et les anciens acteurs internationaux pourrait créer des lacunes dans le partage d’informations et la planification des opérations militaires. Cela affaiblirait l’efficacité globale des efforts antiterroristes dans la région.

Cependant, il convient de noter que ce repositionnement stratégique reflète également une volonté des États sahéliens de diversifier leurs partenariats pour reprendre le contrôle de leur souveraineté. Si cette autonomie stratégique est bien gérée, elle pourrait ouvrir la voie à des solutions mieux adaptées aux besoins locaux.

\quad

\textbf{Conclusion}

La fin de la présence militaire française et le rapprochement avec la Russie marquent un tournant historique pour la région du Sahel. Bien que cette réorientation stratégique puisse offrir de nouvelles opportunités, elle comporte également des risques majeurs pour la stabilité régionale et la lutte contre le terrorisme. Pour éviter une aggravation des crises existantes, il est essentiel que ces changements s'accompagnent d'une gouvernance inclusive, d'un renforcement des coopérations régionales et d'un soutien accru aux initiatives de développement. Le défi pour les États sahéliens sera de maintenir un équilibre entre leurs nouvelles alliances et les aspirations légitimes de leurs populations à la paix et à la sécurité.







\quad

\end{spacing}
\end{document}
